%
% teil1.tex -- Herleitung der Bessel-Differentialgleichung
%
% (c) 2026 Gian Cavegn, OST Ostschweizer Fachhochschule
%
% !TEX root = ../../buch.tex
% !TEX encoding = UTF-8
% !TeX spellcheck = de_CH
%
\section{Herleitung der Bessel-Differentialgleichung
\label{bessel:section:herleitung}}
\kopfrechts{Herleitung}
%
% Ziel dieses Abschnitts.
% Zeigen, dass die Bessel-DGL nicht vom Himmel fällt, sondern
% zwangsläufig entsteht, sobald man ein Wellenproblem auf einem
% Kreisgebiet separiert.
% Roter Faden, Rechteck -> Kreis. Am Rechteck geht alles mit sin/cos
% auf, am Kreis nicht mehr.
%
Die Bessel-Differentialgleichung fällt nicht vom Himmel.
Sie entsteht zwangsläufig, sobald man ein Schwingungsproblem auf einem kreisförmigen Gebiet in Koordinaten zerlegt, die zur Geometrie passen.
Dieser Abschnitt führt die Rechnung von der Wellengleichung bis zur fertigen Differentialgleichung durch.
%
% Die schwingende Membran
%
\subsection{Die schwingende Membran
\label{bessel:herleitung:subsection:membran}}
Eine eingespannte Membran wird durch ihre Auslenkung $u(x,y,t)$ senkrecht zur Ruhelage beschrieben.
Für kleine Auslenkungen genügt $u$ der Wellengleichung
\begin{equation}
\frac{\partial^2 u}{\partial t^2}
=
c^2 \Delta u,
\qquad
\Delta
=
\frac{\partial^2}{\partial x^2}
+
\frac{\partial^2}{\partial y^2},
\label{bessel:herleitung:eqn:welle}
\end{equation}
wobei $c$ die Ausbreitungsgeschwindigkeit ist.
Die Wellengleichung~\eqref{bessel:herleitung:eqn:welle} ergibt sich aus einer Kraftbilanz am Flächenelement.
Die Membran habe die Flächendichte $\mu$ und stehe unter der homogenen Spannung $\sigma$, gemessen als Kraft pro Länge.
An jeder Kante eines Rechtecks mit den Kantenlängen $dx$ und $dy$ zieht die Spannung tangential an der Membran.
Für kleine Auslenkungen ist die vertikale Komponente dieser Kraft proportional zur Steigung, für die beiden Kanten senkrecht zur $x$-Achse also
\[
\sigma\,dy
\bigl(
u_x(x+dx,y) - u_x(x,y)
\bigr)
\approx
\sigma\, u_{xx}\, dx\, dy .
\]
Der Beitrag der beiden anderen Kanten ist $\sigma u_{yy}\,dx\,dy$.
Das zweite Newtonsche Gesetz liefert damit
\[
\mu\, dx\, dy\; u_{tt}
=
\sigma \bigl( u_{xx} + u_{yy} \bigr) dx\, dy ,
\]
und Division durch $\mu\,dx\,dy$ ergibt \eqref{bessel:herleitung:eqn:welle} mit $c^2 = \sigma/\mu$.
Vorausgesetzt sind dabei kleine Auslenkungen, eine über die ganze Membran konstante Spannung und eine verschwindende Biegesteifigkeit.
Eine ausführliche Herleitung findet sich bei \cite{bessel:courant}.
Am Rand der Membran ist die Auslenkung null, es gilt also $u=0$ auf dem Rand des Gebietes.
Für die Zeitabhängigkeit macht man den Separationsansatz
\begin{equation}
u(x,y,t)
=
T(t)\, v(x,y).
\label{bessel:herleitung:eqn:zeitsep}
\end{equation}
Einsetzen in~\eqref{bessel:herleitung:eqn:welle} und Division durch
$c^2Tv$ ergibt
\[
\frac{1}{c^2}\frac{T''(t)}{T(t)}
=
\frac{\Delta v(x,y)}{v(x,y)}.
\]
Die linke Seite hängt nur von $t$ ab, die rechte nur von $x$ und $y$.
Beide Seiten müssen daher konstant sein.
Nennt man diese Konstante $-k^2$, so zerfällt das Problem in die harmonische Schwingungsgleichung $T''=-c^2k^2T$ und in die
\emph{Helmholtz-Gleichung}
\index{Helmholtz-Gleichung}%
\begin{equation}
\Delta v + k^2 v = 0.
\label{bessel:herleitung:eqn:helmholtz}
\end{equation}
Die zeitliche Abhängigkeit ist damit erledigt, sie ist immer sinusförmig mit der Kreisfrequenz $\omega = ck$.
Die ganze Information über die Form der Membran steckt in
\eqref{bessel:herleitung:eqn:helmholtz}.
%
% Rechteck und Kreis
%
\subsection{Rechteck und Kreis
\label{bessel:herleitung:subsection:rechteckkreis}}
An dieser Stelle lohnt sich ein Vergleich, der zeigt, wo die Schwierigkeit liegt.
 
Für eine rechteckige Membran mit den Seitenlängen $a$ und $b$ führt der Ansatz $v(x,y)=X(x)Y(y)$ auf die beiden Gleichungen
$X''=-\alpha^2X$ und $Y''=-\beta^2Y$ mit $\alpha^2+\beta^2=k^2$.
Beide sind Schwingungsgleichungen mit konstanten Koeffizienten, ihre Lösungen sind Sinus- und Kosinusfunktionen, und die Randbedingungen liefern
\[
v_{mn}(x,y)
=
\sin\frac{m\pi x}{a}
\sin\frac{n\pi y}{b},
\qquad
m,n \in \mathbb{N}.
\]
Das Problem ist mit den Funktionen der elementaren Analysis vollständig gelöst.
 
Bei einer kreisförmigen Membran versagt dieser Ansatz.
Der Rand ist ein Kreis, und in kartesischen Koordinaten lässt sich die Randbedingung nicht in eine Bedingung an $X$ und eine an $Y$ zerlegen.
Man muss Koordinaten wählen, in denen der Rand eine Koordinatenfläche ist, und das sind die Polarkoordinaten.
Der Preis dafür ist, dass der Laplace-Operator in diesen Koordinaten nicht mehr konstante Koeffizienten hat.
Genau daraus entsteht die Bessel-Gleichung.
%
% Separation in Polarkoordinaten
%
\subsection{Separation in Polarkoordinaten
\label{bessel:herleitung:subsection:polar}}
In Polarkoordinaten $x=r\cos\varphi$, $y=r\sin\varphi$ lautet der Laplace-Operator
\begin{equation}
\Delta
=
\frac{\partial^2}{\partial r^2}
+
\frac{1}{r}\frac{\partial}{\partial r}
+
\frac{1}{r^2}\frac{\partial^2}{\partial\varphi^2}.
\label{bessel:herleitung:eqn:laplacepolar}
\end{equation}
Diese Form erhält man mit der Kettenregel aus $r=\sqrt{x^2+y^2}$ und $\varphi=\arctan(y/x)$.
Wegen
\[
\frac{\partial r}{\partial x} = \cos\varphi,
\qquad
\frac{\partial r}{\partial y} = \sin\varphi,
\qquad
\frac{\partial\varphi}{\partial x} = -\frac{\sin\varphi}{r},
\qquad
\frac{\partial\varphi}{\partial y} = \frac{\cos\varphi}{r}
\]
lauten die Ableitungsoperatoren
\[
\frac{\partial}{\partial x}
=
\cos\varphi\,\frac{\partial}{\partial r}
-
\frac{\sin\varphi}{r}\,\frac{\partial}{\partial\varphi},
\qquad
\frac{\partial}{\partial y}
=
\sin\varphi\,\frac{\partial}{\partial r}
+
\frac{\cos\varphi}{r}\,\frac{\partial}{\partial\varphi}.
\]
Wendet man beide ein zweites Mal an und addiert, so heben sich die gemischten Ableitungen 
$\partial^2/\partial r\,\partial\varphi$ und die Terme mit $\partial/\partial\varphi$ paarweise weg, 
und mit $\cos^2\varphi+\sin^2\varphi=1$ bleibt\eqref{bessel:herleitung:eqn:laplacepolar} übrig.
Zu beachten ist dabei, dass die Koeffizienten $\cos\varphi$ und $\sin\varphi/r$ selbst von $r$ und $\varphi$ abhängen und deshalb mitdifferenziert werden müssen.
Genau daraus entsteht der Term $\frac{1}{r}\frac{\partial}{\partial r}$.
Die Helmholtz-Gleichung~\eqref{bessel:herleitung:eqn:helmholtz} wird damit zu
\begin{equation}
\frac{\partial^2 v}{\partial r^2}
+
\frac{1}{r}\frac{\partial v}{\partial r}
+
\frac{1}{r^2}\frac{\partial^2 v}{\partial\varphi^2}
+
k^2 v
=
0.
\label{bessel:herleitung:eqn:helmholtzpolar}
\end{equation}
Die Koeffizienten $1/r$ und $1/r^2$ sind der entscheidende Unterschied zum kartesischen Fall.
Sie sind bei $r=0$ singulär, und diese Singularität wird uns im nächsten Abschnitt wieder begegnen.
Nun separiert man ein zweites Mal, diesmal in Radius und Winkel,
\begin{equation}
v(r,\varphi)
=
R(r)\,\Phi(\varphi).
\label{bessel:herleitung:eqn:radialsep}
\end{equation}
Einsetzen in~\eqref{bessel:herleitung:eqn:helmholtzpolar} und Multiplikation mit $r^2/(R\Phi)$ ergibt
\[
\frac{r^2R''(r) + rR'(r) + k^2r^2R(r)}{R(r)}
=
-\frac{\Phi''(\varphi)}{\Phi(\varphi)}.
\]
Wieder hängt die linke Seite nur von $r$ ab und die rechte nur von $\varphi$, beide sind also konstant.
Wir nennen die Konstante $n^2$.
 
Für den Winkelanteil bleibt
\begin{equation}
\Phi''(\varphi) + n^2\Phi(\varphi) = 0
\label{bessel:herleitung:eqn:winkel}
\end{equation}
mit den Lösungen $\Phi(\varphi)=\cos n\varphi$ und $\Phi(\varphi)=\sin n\varphi$.
Hier kommt eine Bedingung ins Spiel, die es beim Rechteck nicht gibt.
Die Winkelkoordinate ist periodisch, es muss $\Phi(\varphi+2\pi)=\Phi(\varphi)$ gelten.
Daraus folgt $n\in\mathbb{Z}$.
Die Separationskonstante ist also nicht frei wählbar, sondern durch die Geometrie auf ganze Zahlen festgelegt.
 
%
% Die Bessel-Differentialgleichung
%
\subsection{Die Bessel-Differentialgleichung
\label{bessel:herleitung:subsection:besseldgl}}
Für den Radialanteil bleibt
\begin{equation}
r^2R''(r) + rR'(r) + (k^2r^2-n^2)R(r) = 0.
\label{bessel:herleitung:eqn:radial}
\end{equation}
Mit der Substitution $x=kr$ und $y(x)=R(x/k)$ verschwindet der
Parameter $k$, denn es ist
$R'(r)=k\,y'(x)$ und $R''(r)=k^2y''(x)$, und man erhält
\begin{equation}
x^2y''(x) + xy'(x) + (x^2-n^2)y(x) = 0.
\label{bessel:herleitung:eqn:bessel}
\end{equation}
Das ist die \emph{Bessel-Differentialgleichung}
\index{Bessel-Differentialgleichung}%
der Ordnung $n$.
 
Aus der Herleitung ergibt sich $n$ als ganze Zahl.
Die Gleichung~\eqref{bessel:herleitung:eqn:bessel} ist aber für jeden reellen oder komplexen Parameter sinnvoll, und die Lösungstheorie des nächsten Abschnitts braucht die Ganzzahligkeit an keiner Stelle.
Wir schreiben deshalb ab jetzt

\begin{equation}
x^2y'' + xy' + (x^2-\nu^2)y = 0,
\qquad
\nu \in \mathbb{R},\ \nu \ge 0,
\label{bessel:herleitung:eqn:besselnu}
\end{equation}
und nennen $\nu$ die \emph{Ordnung}
\index{Ordnung einer Bessel-Funktion}%
der Gleichung.
Halbzahlige Ordnungen treten zum Beispiel bei kugelsymmetrischen Problemen auf, wo dieselbe Rechnung in Kugelkoordinaten durchgeführt wird.
 
Ein Blick auf~\eqref{bessel:herleitung:eqn:besselnu} zeigt schon, wo die Schwierigkeit liegt.
Dividiert man durch $x^2$, so lautet die Gleichung
\[
y''
+
\frac{1}{x}y'
+
\frac{x^2-\nu^2}{x^2}y
=
0,
\]
und die Koeffizienten haben bei $x=0$ einen Pol.
Der Existenzsatz für Potenzreihenlösungen aus Kapitel~\ref{chapter:differentialgleichungen} ist daher nicht anwendbar.
Tatsächlich ergibt der Ansatz $y=\sum_k a_kx^k$ die Bedingung $(k^2-\nu^2)a_k=-a_{k-2}$, und für nicht ganzzahliges $\nu$ bleibt davon nur die triviale Lösung übrig.
Die Gleichung ist aber von genau der Bauart
\[
x^2y'' + p(x)\,xy' + q(x)\,y = 0
\]
mit den bei $0$ analytischen Funktionen $p(x)=1$ und $q(x)=x^2-\nu^2$, sie ist also eine fuchssche Differentialgleichung im Sinne von Abschnitt~\ref{buch:differentialgleichungen:section:fuchs}.
Damit steht das Werkzeug für ihre Lösung bereits zur Verfügung.